\chapter{Problem Definition and Formulation}
\label{ch:problem}

This chapter provides a rigorous mathematical definition of the SOH estimation problem and formulates the specific physical models that govern the hybrid framework.

\section{Problem Definition}

The objective of the thesis is to estimate the State of Health ($SOH$) of a Lithium-ion battery at a given cycle index $k$, based on a sequence of historical measurements. State of Health (SOH) is defined as the ratio of the maximum releasable capacity at the current cycle ($Q_{max, k}$) to the nominal rated capacity ($Q_{rated}$) provided by the manufacturer:

\begin{equation}
    SOH_k = \frac{Q_{max, k}}{Q_{rated}} \times 100\%
\end{equation}

The problem is modeled as a supervised time-series regression task. Let $\mathbf{x}_k \in \mathbb{R}^d$ be a vector of $d$ features extracted at cycle $k$ (e.g., voltage, current, ICA peak position, temperature). The prognostic model $f(\cdot)$ aims to map a sequence of historical features of length $T$ to the current SOH:

\begin{equation}
    \hat{y}_k = f(\mathbf{X}_{k-T:k}; \theta)
\end{equation}

Where $\mathbf{X}_{k-T:k}$ represents the input tensor of historical data, $\theta$ represents the trainable parameters of the neural network (GRU), and $\hat{y}_k$ is the predicted SOH.

\textbf{End-of-Life (EOL) Definition:} End-of-Life (EOL) is defined as the first cycle index where $SOH \leq SOH_{EOL}$ \emph{and the condition persists for at least 5 consecutive cycles} (this persistence rule prevents a transient noisy reference cycle or a regeneration dip from triggering a false EOL). The primary threshold is $SOH_{EOL} = 80\%$ for both datasets (matching common automotive practice); a $70\%$ threshold is sometimes adopted for less demanding second-life applications, but all results in this thesis use the $80\%$ threshold. Both datasets are LCO, so the threshold choice reflects the application rather than a chemistry difference. All RUL calculations explicitly document the EOL threshold used.

\textbf{RUL Definition:} Given the EOL cycle index $k_{EOL}$ determined by the rule above, the Remaining Useful Life at cycle $k$ is the number of cycles left until that point:
\begin{equation}
    RUL_k = k_{EOL} - k, \qquad k \leq k_{EOL}.
\end{equation}
$RUL_k$ is undefined for $k > k_{EOL}$ and for cells whose record never reaches the threshold (B0007 never crosses even the 70\% line within its record). Because $k_{EOL}$ depends on the chosen threshold and persistence rule, every reported RUL value carries that definition with it.

\section{Physics Formulation (The Regularizer)}
\label{sec:physics_formulation}

To ensure the neural network adheres to physical laws, we introduce a regularization term based on empirical degradation physics. We employ two candidate Ordinary Differential Equations (ODEs) as the governing physical laws. These ODEs act as ``soft constraints'' on the loss function.

\subsection{The Verhulst (Logistic) Model}

The Verhulst model describes a system whose growth (or decay) is self-limiting. In the context of batteries, it implies that the degradation rate is proportional to the current state and the remaining ``capacity'' for degradation (saturation of SEI layer growth). The governing differential equation is:

\begin{equation}
    \frac{dSOH}{dt} = r \cdot SOH \cdot \left( 1 - \frac{SOH}{K} \right)
\end{equation}

Where:
\begin{itemize}
    \item $r$: The intrinsic rate constant. With $r>0$ the logistic law describes saturating growth toward $K$; fitted to a \emph{decaying} SOH trajectory (or when learned jointly with the network, Chapter~\ref{ch:results}), $r$ takes a negative sign and the same law describes monotone fade.
    \item $K$: The carrying capacity. For $r>0$ it is the asymptotic plateau of the S-curve; fitted with $r<0$ to a \emph{decaying} SOH trajectory (the case throughout Chapter~\ref{ch:results}), the closed-form curve $K/(1+A e^{-rt})$ instead decays from its starting value $K/(1+A)$ toward zero, so $K$ acts as a scale parameter of the fitted curve rather than an asymptote the trajectory approaches. Because SOH is normalized by \emph{rated} capacity, fitted values of $K$ may exceed $1.0$ for cells whose measured initial capacity is above rating (as observed for the NASA cells in Chapter~\ref{ch:results}).
\end{itemize}

This model enforces a specific trajectory shape: an initial slow fade, a near-linear phase, and a saturation phase. By constraining the network's time-derivative to this law, we discourage non-physical predictions such as sustained capacity recovery ($\frac{dSOH}{dt} > 0$), which would violate the monotonic degradation logic inherent to thermodynamic aging \cite{ref1_general}.

\subsection{The Double-Exponential Model}

To capture more complex degradation dynamics, particularly the ``knee-point'' where aging accelerates (often due to lithium plating or massive loss of active material), we formulate a Double-Exponential prior. This model assumes the total degradation is the superposition of two independent exponential decay processes:

\begin{equation}
    SOH(t) = a \cdot e^{b \cdot t} + c \cdot e^{d \cdot t}
\end{equation}

\begin{itemize}
    \item \textbf{Term 1 ($a \cdot e^{b \cdot t}$):} Represents the fast decay process, typically dominated by initial SEI formation and stabilization in the early cycles.
    \item \textbf{Term 2 ($c \cdot e^{d \cdot t}$):} Represents the slow decay process, dominated by the long-term Loss of Active Material (LAM).
\end{itemize}

This mechanistic reading motivates the functional form but is not imposed on
the least-squares fit: the coefficients are left unconstrained, and, as noted
in Chapter~\ref{ch:results}, the fitted amplitudes and exponents do not always
take the signs of two decay processes. The double-exponential is therefore
used as a flexible empirical prior whose \emph{form} is physically motivated,
rather than as a literal decomposition into SEI and LAM contributions.

For the PINN framework, we differentiate this equation to obtain the rate of change constraint:

\begin{equation}
    \frac{dSOH}{dt} = a \cdot b \cdot e^{b \cdot t} + c \cdot d \cdot e^{d \cdot t}
\end{equation}

This model offers greater flexibility than the Verhulst model, allowing the network to fit datasets with sharp transitions in degradation rates, directly addressing the limitations of simple logistic models \cite{wang_pinn}.

\textbf{Soft Constraints \& Handling Regeneration:} The ODE priors (Verhulst, Double-Exponential) are implemented as soft regularizers rather than hard constraints. The uncertainty-weighted loss allows the training procedure to attenuate the influence of the physics loss when transient phenomena (e.g., capacity regeneration in NASA data) cannot be explained by the ODE prior. This avoids forced suppression of legitimate data observations while still guiding the model toward physically plausible trends.

\section{The Optimization Objective}
\label{sec:optimization_objective}

The core innovation of this framework lies in the composite loss function. Unlike traditional deep learning, which minimizes only the error between prediction and ground truth, our framework minimizes a joint objective. This composite loss function ensures the model fits the data while respecting the physical laws governing battery aging.

\begin{equation}
    \mathcal{L}_{total} = \lambda_{data}\mathcal{L}_{data} + \lambda_{phys}\mathcal{L}_{phys}
    \label{eq:composite_loss}
\end{equation}


\textbf{Data Loss ($\mathcal{L}_{data}$):} We use the Mean Squared Error (MSE) to measure the discrepancy between the network's prediction $\hat{y}$ and the ground truth labels $y$:
\begin{equation}
    \mathcal{L}_{data} = \frac{1}{N} \sum_{i=1}^{N} (y_i - \hat{y}_i)^2
\end{equation}

\textbf{Physics Loss ($\mathcal{L}_{phys}$):} We use Automatic Differentiation (Autograd) to compute the time derivative of the network's output ($\frac{d\hat{y}}{dt}$) with respect to the input time/cycle index. Strictly, this is the \emph{partial} derivative of the network output with respect to its cycle-coordinate input, holding the other input features fixed; it serves as a tractable proxy for the trajectory derivative (the remaining features are not differentiable functions of the cycle index), which is standard practice in physics-informed training. We then compare this computed derivative to the expected derivative derived from the ODE ($g(\hat{y}, \theta_{phys})$):
\begin{equation}
    \mathcal{L}_{phys} = \frac{1}{N} \sum_{i=1}^{N} \left| \frac{d\hat{y}}{dt} - g(\hat{y}, \theta_{phys}) \right|^2
\end{equation}

This formulation ensures that the model fits the training data while ensuring the \textbf{trajectory} of the prediction is consistent with the underlying physics of degradation.

The ODE parameters $\theta_{phys}$ admit two treatments, and the framework implements both. They may be fit once on the training cells and held fixed, so that the physics term penalizes deviation from a known closed-form curve (route~A in Chapter~\ref{ch:methodology}); or they may be implemented as learnable parameters within the computation graph and updated via backpropagation alongside the network weights, allowing the physical law to adapt its shape to the training data (route~B). Chapter~\ref{ch:results} compares the two treatments empirically rather than assuming either is superior.

\par\medskip
\noindent \textbf{Note on Implementation Strategy:} 
Equation~(\ref{eq:composite_loss}) represents the canonical form of a multi-objective physics-informed loss function, where the balance between data fidelity and physical consistency is controlled by static scalar weights ($\lambda_{data}, \lambda_{phys}$). Determining suitable values for these weights is non-trivial. Chapter~\ref{ch:methodology} (Section~\ref{sec:uncertainty_weighting}) therefore develops, in addition to the fixed-weight form above, a probabilistic alternative in which the static hyperparameters $\lambda$ are replaced by learnable uncertainty terms ($\sigma$) intended to adapt the balance automatically. Both weighting schemes are evaluated empirically in Chapter~\ref{ch:results}; on the datasets studied there, the learnable weighting brought no reliable accuracy gain over a well-chosen fixed weight, which is therefore used for the main benchmark models.